% quant-lab / dart-event-study 종합 결과 리포트
% 컴파일: xelatex quant_lab_report.tex   (한글 폰트 필요 — Windows 기본 '맑은 고딕')
% 모든 수치는 data/ 산출물(parquet/json)에서 직접 추출. 임의 반올림·선별 없음.
\documentclass[11pt,a4paper]{article}

\usepackage{kotex}
\usepackage[margin=2.4cm]{geometry}
\usepackage{booktabs}
\usepackage{amsmath}
\usepackage{longtable}
\usepackage{array}
\usepackage{xcolor}
\usepackage[colorlinks=true,linkcolor=black,urlcolor=blue!60!black]{hyperref}
\usepackage{fancyhdr}

\pagestyle{fancy}
\fancyhf{}
\fancyhead[L]{\small quant-lab · DART 이벤트 스터디}
\fancyhead[R]{\small\thepage}
\renewcommand{\headrulewidth}{0.4pt}

\newcommand{\sig}[1]{\textbf{#1}}
\newcommand{\ns}{\textcolor{gray}{n.s.}}

\title{\vspace{-1.5cm}\textbf{한국 주식 공시 이벤트의 정보력과 체결가능성}\\
\large 이벤트 스터디 · 검정 배터리 · 리테일 어텐션 상호작용}
\author{quant-lab / \texttt{dart-event-study}}
\date{2026년 7월 26일}

\begin{document}
\maketitle
\thispagestyle{fancy}

\begin{abstract}
\noindent
KOSPI 시총 상위 200 종목(2019-01 기준, 상장폐지 포함)의 2019--2024년 DART 공시에서
자사주 취득·유상증자·잠정실적 이벤트를 추출해 시장모형 이벤트 스터디와 백테스트를 수행했다.
핵심 결과는 세 가지다. 첫째, \sig{자사주 취득 공시 후 20거래일 누적초과수익(CAAR) $+4.88\%$}는
BMP·Corrado 순위·부호검정·부트스트랩·플라시보 순열검정을 모두 통과하는 유일한 강건한 신호다.
둘째, 이 정보력은 \sig{체결가능한 초과수익으로 이어지지 않는다} --- 익영업일 종가 진입과 현실적
거래비용 아래 연 $+2.4\%$(KOSPI $+3.1\%$)로 벤치마크를 밑돌고, Deflated Sharpe $3.5\%$로
우연과 구분되지 않는다. 셋째, 뉴스 어텐션과 드리프트의 \sig{상호작용은 존재하지 않는다}(전 스펙
$p \ge 0.22$). 플라시보 검정은 실적 공시의 겉보기 장기 드리프트가 위양성임을 보였고, 반대로
저검정력 모수검정이 놓친 유상증자 효과를 되살렸다. 본 리포트는 유의한 결과만이 아니라
기각·귀무 결과를 동일한 비중으로 기록한다.
\end{abstract}

\section{데이터와 방법}

\subsection{표본}
유니버스는 2019-01 프록시 시가총액(발행주식총수 $\times$ 2019-01 수정종가) 상위 200 보통주이며,
기간 내 상장폐지 11종목을 포함해 생존편향을 보정했다. 공시 139{,}088건에서 이벤트 5{,}078건
(자사주 직접 178 / 신탁 106 / 실적 5{,}300 / 유상증자 139)을 추출하고 방향성 시그널 4{,}650건을 얻었다.
가격은 FinanceDataReader(네이버) 수정주가를 쓴다. KRX 정보데이터시스템은 봇 차단으로 사용 불가임을
실측 확인했다.

\subsection{이벤트 스터디 설계}
정상수익률은 시장모형 $R_{i,t} = \alpha_i + \beta_i R_{m,t} + \varepsilon_{i,t}$로 추정한다.
추정창은 이벤트 이전 120거래일이며 이벤트와 10거래일의 갭을 둔다. 이벤트일(day 0)은 접수일 이후
첫 거래일이다. 누적초과수익 $\mathrm{CAR}_i[a,b] = \sum_{t=a}^{b} \mathrm{AR}_{i,t}$를
사전 지정한 세 윈도우 $[-1,+1]$, $[0,+5]$, $[0,+20]$에서 계산하고 \emph{전부} 보고한다.

\subsection{Look-ahead 통제}
접수시각은 OpenDART가 제공하지 않음을 실측 확인했으므로(DART 뷰어·KIND도 과거 시각 미제공),
전 건에 \sig{익영업일 종가 체결}이라는 보수적 가정을 적용했다. 이는 당일 체결 기회를 포기하는
방향의 편향이므로 look-ahead가 아니다. 실적 서프라이즈는 외부 컨센서스 대신 공시 원문의
전년동기 수치로 YoY를 계산해 look-ahead를 원천 차단했다. 발행주식총수는 이벤트일 이전에
법정기한상 공개된 정기보고서 기준으로만 참조한다.

\subsection{이벤트 클러스터링}
실적 공시는 어닝시즌에 날짜가 뭉친다(시그널 4{,}294건이 523일에 발생, 최대 하루 48건).
이벤트 간 독립을 가정하는 naive $t$는 이 경우 과대추정되므로, 상수항 회귀의 클러스터-로버스트
분산(CR1, 자유도 $G-1$)을 월 단위 클러스터로 계산해 병기한다. 모든 유의성 판정은 월클러스터
기준이며, 15개 검정(5그룹 $\times$ 3윈도우)에 Benjamini--Hochberg FDR 보정을 적용한다.

\section{결과 1: 공시의 정보력}

\begin{table}[htbp]
\centering
\caption{그룹별 CAAR과 유의성. $t$는 naive $\to$ 월클러스터. $q$는 BH-FDR 보정값.}
\small
\begin{tabular}{llrrrrr}
\toprule
그룹 & 윈도우 & $N$ & CAAR & $t_{\text{naive}}$ & $t_{\text{cl}}$ & $q$ \\
\midrule
자사주 직접 & $[-1,+1]$ & 169 & $+3.05\%$ & 8.25 & 7.47 & \sig{$<0.001$} \\
             & $[0,+5]$  & 169 & $+3.84\%$ & 7.97 & 6.44 & \sig{$<0.001$} \\
             & $[0,+20]$ & 169 & \sig{$+4.88\%$} & 5.82 & 4.88 & \sig{$<0.001$} \\
\addlinespace
자사주 신탁 & $[-1,+1]$ & 104 & $+2.56\%$ & 5.22 & 4.62 & \sig{$0.0001$} \\
             & $[0,+5]$  & 104 & $+2.19\%$ & 2.79 & 2.16 & $0.060$ \ns \\
             & $[0,+20]$ & 104 & \sig{$+4.10\%$} & 3.13 & 2.66 & \sig{$0.027$} \\
\addlinespace
실적 $+$ & $[-1,+1]$ & 2391 & $+1.26\%$ & 11.61 & 6.90 & \sig{$<0.001$} \\
         & $[0,+5]$  & 2391 & $+0.53\%$ & 3.83 & 2.24 & $0.059$ \ns \\
         & $[0,+20]$ & 2391 & $+0.10\%$ & 0.43 & 0.20 & $0.845$ \ns \\
\addlinespace
실적 $-$ & $[-1,+1]$ & 1767 & $-0.53\%$ & $-4.45$ & $-2.61$ & \sig{$0.027$} \\
         & $[0,+5]$  & 1767 & $-0.71\%$ & $-4.49$ & $-2.20$ & $0.059$ \ns \\
         & $[0,+20]$ & 1767 & $-0.18\%$ & $-0.73$ & $-0.31$ & $0.801$ \ns \\
\addlinespace
유상증자 $-$ & $[-1,+1]$ & 49 & $-1.88\%$ & $-1.74$ & $-1.76$ & $0.123$ \ns \\
             & $[0,+5]$  & 49 & $-3.93\%$ & $-3.28$ & $-2.93$ & \sig{$0.019$} \\
             & $[0,+20]$ & 49 & $-4.98\%$ & $-2.09$ & $-1.79$ & $0.123$ \ns \\
\bottomrule
\end{tabular}
\end{table}

FDR 보정 후 살아남는 것은 자사주(직접·신탁)의 전 윈도우 드리프트, 실적 $\pm$의 공시창
$[-1,+1]$, 유상증자 $-$의 $[0,+5]$다. 경계 유의였던 실적 $[0,+5]$($p\approx0.03$)는 다중검정
보정에서 탈락했으며, 탈락한 대로 기록한다. 유상증자 $+$($N=22$)는 표본이 얇아 결론을 유보한다.

\section{결과 2: 검정 배터리와 플라시보}

단일 $t$검정에 의존하지 않기 위해 네 가지 추가 검정을 수행했다.
\sig{BMP}(Boehmer--Musumeci--Poulsen)는 각 CAR를 자기 추정기간의 예측오차 표준편차로 표준화해
이벤트유발 분산(event-induced variance)을 보정한다. \sig{Corrado 순위검정}과 \sig{일반화 부호검정}
(Cowan, 귀무 양수확률을 추정기간 실제 비율로 설정)은 정규성·등분산 가정을 제거한다.
\sig{부트스트랩}은 이벤트 단위 복원추출로 CAAR의 신뢰구간을 구한다.

가장 중요한 것은 \sig{플라시보 순열검정}이다. 각 이벤트를 \emph{같은 종목의 무작위 거래일}로
치환하되 실제 이벤트 $\pm30$거래일은 제외하고, 동일 파이프라인을 1{,}000회 반복해 CAAR의
귀무분포를 만든다. 종목 구성과 표본 크기를 고정한 채 타이밍만 무작위화하므로,
``이 효과가 아무 날에나 나오는가''를 직접 답한다.

\begin{table}[htbp]
\centering
\caption{검정 배터리 요약(윈도우 $[0,+20]$ 및 주요 창). 플라시보 $p$의 해상도 하한은 $1/1001$.}
\small
\begin{tabular}{llrrrrr}
\toprule
그룹·윈도우 & CAAR & BMP $p$ & Corrado $p$ & 부호 $p$ & 플라시보 귀무평균 & 플라시보 $q$ \\
\midrule
자사주 직접 $[0,+20]$ & $+4.88\%$ & \sig{$<0.001$} & \sig{$<0.001$} & \sig{$<0.001$} & $-0.21\%$ & \sig{$0.002$} \\
자사주 신탁 $[0,+20]$ & $+4.10\%$ & \sig{$0.0006$} & \sig{$0.002$} & \sig{$0.023$} & $-0.28\%$ & \sig{$0.002$} \\
실적 $+$ $[0,+20]$ & $+0.10\%$ & $0.002$ & $0.970$ & $0.001$ & $+0.11\%$ & $0.633$ \ns \\
실적 $-$ $[0,+20]$ & $-0.18\%$ & $0.662$ & $0.462$ & $0.646$ & $-0.34\%$ & $0.840$ \ns \\
유상증자 $-$ $[0,+5]$ & $-3.93\%$ & $0.143$ \ns & \sig{$0.002$} & \sig{$0.0007$} & $-0.02\%$ & \sig{$0.011$} \\
\bottomrule
\end{tabular}
\end{table}

플라시보 검정은 세 가지를 밝혔다.

\begin{enumerate}
\item \sig{자사주 드리프트는 확증된다.} 무작위 날짜의 귀무 CAAR 평균은 $-0.21\%$(직접),
$-0.28\%$(신탁)로 0 근처이고, 실제 $+4.88\%$/$+4.10\%$는 경험적 $p$가 해상도 하한에 닿는다.
헤드라인 결과가 가장 엄격한 sanity check를 통과한다.

\item \sig{실적 $+$의 장기 드리프트는 위양성이다.} BMP($p=0.002$)와 부호검정($p=0.001$)은
유의해 보이지만, 이는 귀무가설을 CAAR${}=0$으로 놓기 때문이다. 대형 성장주는 임의의 20거래일
구간에서도 조금 오른다. 플라시보 귀무평균 $+0.11\%$는 실제값 $+0.10\%$와 사실상 동일하며
($p=0.56$, Corrado $p=0.97$), 따라서 이는 \emph{이벤트 효과가 아니라 기저 드리프트}다.
\emph{CAAR${}=0$이라는 귀무가설 자체가 잘못 교정되어 있었다}는 것이 이 검정의 기여다.

\item \sig{유상증자 효과는 되살아난다.} $N=49$의 두꺼운 꼬리 탓에 모수검정 BMP는 비유의
($p=0.14$)지만, Corrado($p=0.002$)·부호($p=0.0007$)·부트스트랩(95\% CI $[-6.09\%, -1.52\%]$)·
플라시보($q=0.011$)가 일관되게 유의하다. 희석 악재는 실재한다.
\end{enumerate}

\section{결과 3: 체결가능성}

이벤트 스터디의 정보력이 거래 가능한 초과수익인지 확인하기 위해, 익영업일 종가 진입·동시
포지션 균등가중·$|\Delta w|$ 기반 비용 부과로 백테스트했다. 거래세는 시행일별 스케줄
(2019년 0.30\% $\to$ 2024년 0.18\%, 매도측)에 슬리피지 0.10\%와 위탁수수료 0.015\%(편도)를 더한다.

\begin{table}[htbp]
\centering
\caption{보정 사다리. 각 단계는 편향을 하나씩 제거한다.}
\small
\begin{tabular}{lrr}
\toprule
보정 단계 & net 연수익 & Sharpe \\
\midrule
v1.1: 생존자 스냅샷 유니버스 $+$ 고정 세율 0.20\% & $+5.1\%$ & 0.21 \\
$+$ 2019년 기준 유니버스·상폐 처리(승자편향 제거) & $+3.8\%$ & 0.17 \\
$+$ 시행일별 실효 세율 $+$ 수수료 & \sig{$+2.4\%$} & \sig{0.10} \\
\midrule
KOSPI 벤치마크 & $+3.1\%$ & 0.16 \\
\bottomrule
\end{tabular}
\end{table}

보정을 마친 자사주 $H=5$ 전략은 연 $+2.4\%$, Sharpe 0.10으로 \sig{벤치마크(0.16)를 밑돈다}.
블록 부트스트랩(블록 20일, 2{,}000회) 95\% 신뢰구간은 $[-0.57, +0.97]$이고 $P(\text{Sharpe} \le 0) = 29.5\%$다.
시도한 20개 조합을 반영한 \sig{Deflated Sharpe는 3.5\%}로, 관측 Sharpe 0.218이 우연히 기대되는
최대치 0.961에 크게 못 미친다.\footnote{프로젝트 README는 이 값을 7\%로 기록했으나, 현재
산출물 \texttt{significance\_full.json}의 값은 3.5\%다. 백테스트 재실행 과정에서 시도 조합
집합이 갱신된 결과로 보이며, 본 리포트는 산출물 값을 따른다. 두 값 모두 ``우연과 구분 불가''라는
결론은 동일하다.}

서브기간 분해는 성과가 국면에 의존함을 보인다. 자사주 $H=5$는 2019--21년 연 $-7.2\%$
(KOSPI $+14.3\%$)인 반면 2022--24년 $+16.3\%$(KOSPI $-7.3\%$)로, 밸류업 담론기에 집중된다.
실적 시그널은 회전율이 연 43--55회에 달해 실효 비용 반영 후 연 $-20\%$로 거래 불가가 확정된다.

민감도는 취약하다. 편도 슬리피지를 0.1\%/0.3\%/0.5\%로 바꾸면 자사주 $H=5$의 net 연수익은
$+2.4\%$/$-7.5\%$/$-16.5\%$로 반전된다. 실효 비용이 조금만 높아도 결론이 뒤집히는 마진이다.

\section{결과 4: 리테일 어텐션 상호작용}

``공시 후 드리프트가 리테일 어텐션 수준에 따라 증폭되는가''를 검증했다. 어텐션은 접수일
$[-1,+1]$ 창의 네이버 뉴스 기사 수이며, 자사주 284건과 유상증자 139건에 대해 수집했다
(실적 5{,}300건은 스로틀 기준 약 10시간이 소요되고 차단 위험이 커 제외).

\subsection{설계상의 함정과 사전 등록}
어텐션 수집창이 $[-1,+1]$이므로 이 변수는 \sig{사전 관심이 아니라 사후·동시 반응}이다.
따라서 어텐션 창과 겹치는 $\mathrm{CAR}[0,+20]$과 겹치지 않는 $\mathrm{CAR}[+2,+20]$을
사전에 모두 등록하고, 판정 규칙도 미리 고정했다: \emph{겹칠 때만 유의하고 겹치지 않을 때
무의미하면 그것은 예측력이 아니라 동시성이다.}

회귀식은 event\_type을 더미로 일반화한 형태다(기준 범주 = 자사주 직접).
\[
\mathrm{CAR}_i = \beta_0 + \sum_k \gamma_k D_{k,i} + \delta\,\mathrm{attn}_i
+ \sum_k \theta_k (D_{k,i} \times \mathrm{attn}_i) + \mathbf{x}_i'\boldsymbol{\lambda} + u_i
\]
$\mathrm{attn}$은 $\log(1+\text{기사 수})$의 표본 z-score이고, 통제변수는 로그 시가총액과
추정창 잔차 변동성이다. 표준오차는 월 클러스터-로버스트다.

\begin{table}[htbp]
\centering
\caption{상호작용 회귀. 주요 스펙 = $\mathrm{CAR}[0,+20]$, 통제 포함($N=357$, 클러스터 67, $R^2=0.131$).}
\small
\begin{tabular}{lrrr}
\toprule
항 & 계수 & 표준오차 & $p$ \\
\midrule
$\mathrm{attn}$ (기준=자사주 어텐션 기울기) & $+0.0197$ & 0.0103 & $0.061$ \\
$\mathrm{attn} \times$ 신탁 \quad(\textbf{상호작용}) & $+0.0009$ & 0.0204 & $0.963$ \ns \\
$\mathrm{attn} \times$ 유상증자 \quad(\textbf{상호작용}) & $-0.0209$ & 0.0167 & $0.216$ \ns \\
신탁 더미 & $+0.0107$ & 0.0176 & $0.544$ \ns \\
유상증자 더미 & $-0.0650$ & 0.0153 & \sig{$<0.001$} \\
로그 시가총액 & $-0.0079$ & 0.0051 & $0.130$ \\
추정창 변동성 & $-2.2807$ & 0.8953 & \sig{$0.013$} \\
\bottomrule
\end{tabular}
\end{table}

\sig{상호작용은 네 개 스펙 전부에서 귀무다.} $\mathrm{attn}\times$신탁은 $p \ge 0.70$,
$\mathrm{attn}\times$유상증자는 $p \ge 0.22$이며 부호도 스펙에 따라 바뀐다. 방향이 반대인
유상증자(희석 악재)를 포함해도 어텐션 효과가 이벤트 타입별로 다르다는 증거는 없다.

어텐션 주효과 역시 예측력으로 보기 어렵다. 유의에 근접한 것은 어텐션 창과 겹치는 스펙
하나뿐이고($p=0.061$), 겹치지 않는 $[+2,+20]$에서는 사라진다($p=0.28$). 사전에 고정한 규칙에
따라 이는 \sig{동시성이지 예측력이 아니다}. $R^2$는 0.036--0.131에 그친다.

한편 유상증자 더미가 $-6.5\%p$로 강하게 유의한 것은 모델이 알려진 희석 악재 드리프트를
정확히 포착한다는 검증 신호다. 부수적으로, 자사주만으로 회귀했을 때 유의했던 로그 시가총액
($p=0.002$)이 유상증자 추가 후 약해지고($p=0.13$) 추정창 변동성이 유의해진다. 삼분위 정렬에서
관찰됐던 ``관심이 많을수록 드리프트가 크다''는 단조 패턴의 상당 부분은 \sig{규모 교락}이었다.

\section{한계}

\begin{itemize}
\item \textbf{유니버스 프록시.} 2019-01 시총을 ``현재 주식수 $\times$ 2019-01 수정종가''로 근사했다.
이후 유상증자를 한 종목은 2019년 시총이 과대평가되며, 기간 중 신규 상장 종목은 미편입이다.
실제 KOSPI200 구성 이력과는 다르다.
\item \textbf{접수시각 미확보.} 전 건 익영업일 체결 가정은 보수적 편향이다. 접수시각 소스를
확보하면 당일 체결 변형이 가능하다(로직은 구현·테스트 완료).
\item \textbf{실적 이벤트 약 13\%는 방향 판정 불가}(2019--20년 원문 미제공 + 표 형식 변형).
\item \textbf{컨센서스 부재.} 서프라이즈는 YoY 대용치이며 애널리스트 기대 대비가 아니다.
\item \textbf{배당 제외} 가격수익률이다(벤치마크 KOSPI 가격지수와 일관).
\item \textbf{어텐션 표본의 편중.} 뉴스 어텐션은 자사주·유상증자에만 수집됐고 실적은 제외됐다.
유상증자 방향은 혼합(제3자배정 65 / 악재 배정 50 / 호재 24)이다.
\item \textbf{숏 제약 미반영.} 유상증자 숏 전략의 대차 가능성·비용을 모델링하지 않았다.
\end{itemize}

\section{부록: 리테일 커뮤니티 수집 현황}

두 번째 축으로 토스증권 종목토론 커뮤니티를 순방향 수집했다(2026-07-15 개시, 19종목,
31{,}971글, poll 61회). 토스는 절대 시각을 노출하지 않으므로 글이 처음 관측된 크롤 시각을
\texttt{first\_seen}으로 기록하는 설계를 택했다.

수집 데이터의 기술통계는 다음과 같다. 종목별 활동량은 일 60--334글로 팍스넷(거래일당 약 33글)
대비 밀도 우위가 확인된다. 장중/장후 분리는 작동한다(장중 22{,}601 / 장후 7{,}006 / 주말 2{,}364,
장중 비중 70.7\%).

그러나 \sig{삭제 이벤트 관측은 현재 스케줄에서 사실상 불가능하다}. 글당 재관측은 평균 1.08회이며
\sig{94.7\%가 단 1회만 관측}된다. 삭제 판정은 대상 글의 피드상 위·아래 이웃이 이후 poll에서
모두 재관측되어야 성립하는데(전역 id 경계를 쓰면 고정글 아웃라이어 때문에 대량 오탐이 발생한다 ---
실측 000660에서 736건), 하루 4슬롯으로는 유입량을 따라가지 못해 대부분의 글이 다음 크롤 전에
관측 깊이 밖으로 밀려난다. 결과적으로 삭제 확정은 39건이고 \sig{97.9\%가 판정 불가(censored)}다.
여기서 censored는 ``삭제되지 않았다''가 아니라 ``판정할 수 없다''는 뜻이며, 이를 삭제 0으로
집계하지 않는 것이 설계의 핵심이다.

이 축은 아직 \sig{DART 이벤트와 시간이 겹치지 않는다}(토스 2026-07~ vs 이벤트 2019--2024).
따라서 어텐션-드리프트 접합 분석은 수개월의 축적 이후에나 가능하며, 본 리포트의 어텐션 분석은
전적으로 뉴스 데이터에 기반한다.

\section{결론}

공시의 정보력은 실재하며 유니버스·통계 보정을 모두 통과한다. 자사주 취득 공시의
$[0,+20]$ 드리프트 $+4.88\%$는 모수·비모수·플라시보 검정을 일관되게 통과하는 유일한 강건
신호다. 그러나 v1.1에서 ``간신히 생존''으로 보였던 체결가능 초과수익은 상당 부분 승자 편향과
비용 과소평가의 산물이었고, 보정 후에는 벤치마크를 밑돌며 우연과 구분되지 않는다
(Deflated Sharpe 3.5\%).

리테일 어텐션이 드리프트를 조절한다는 가설은 지지되지 않았다. 상호작용 계수는 모든 스펙에서
귀무이고, 어텐션 주효과로 보였던 것은 어텐션 측정창과 수익률 측정창이 겹치는 데서 오는 동시성이었다.

이 프로젝트의 기여는 수익률이 아니라 \sig{이 괴리의 정량화}와 \sig{기각 과정의 기록}이다.
특히 플라시보 검정이 실적 공시의 겉보기 유의를 걷어내고 유상증자 효과를 되살린 사례는,
$\mathrm{CAAR}=0$이라는 관행적 귀무가설이 드리프트하는 자산에 대해 잘못 교정되어 있음을 보여준다.

\vspace{0.5cm}
\noindent\rule{\textwidth}{0.4pt}\\
\small
\textbf{재현.} 전 파이프라인은 \texttt{python -m dart\_event\_study <stage>}로 실행한다
(\texttt{--list}로 전체 단계 확인). 유닛테스트 162개는 네트워크·API 키 없이 실행되며,
GitHub Actions가 push마다 Ubuntu·Windows 양쪽에서 ruff 린트와 pytest를 수행한다.
모든 난수 시드는 고정돼 있다. 본 리포트의 수치는 \texttt{data/} 산출물에서 직접 추출했다.

\end{document}
